\chapter{Complexity — Uniform Tableau Emitter Layer}
%%% generated by the blueprint pipeline from TCSlib.Complexity.CircuitComplexity.UniformTableauEmitterLayer
\section{Overview}

Declarations of the Lean module \texttt{TCSlib.Complexity.CircuitComplexity.UniformTableauEmitterLayer}.

\section{Declarations}

\begin{definition}[Instruction template bound]
\lean{Complexity.UTab.Ki}
\label{Complexity.UTab.Ki}
\leanok
\uses{Complexity.UTab.Lmax}
For a finite binary Turing machine $M$, the instruction template bound is
$L_{\max}(M)+3$, where $L_{\max}(M)$ is the maximum length of its instruction templates.
\end{definition}

\begin{theorem}[Execution of a template]
\lean{Complexity.UTab.goes_instr_inv}
\label{Complexity.UTab.goes_instr_inv}
\leanok
\uses{Complexity.CounterProg.Goes, Complexity.CounterProg.Goes.congr, Complexity.CounterProg.MOp.exec, Complexity.UTab.ER, Complexity.UTab.ER.f, Complexity.UTab.Ki, Complexity.UTab.Lb, Complexity.UTab.Lmax, Complexity.UTab.Tpl, Complexity.UTab.Tpl.after, Complexity.UTab.Tpl.next, Complexity.UTab.Tpl.t0, Complexity.UTab.goes_instr, Complexity.UTab.length_tmpl_le, Complexity.UTab.prog, Complexity.UTab.tmpl, Complexity.UTab.tpL}
\difficulty{4}
Let $M$ be a finite Turing machine with Boolean symbols, let $x$ be a Boolean word, and
let $k$ be the number of its work tapes. For a template $\theta$ for $M$, suppose its
successor label is the label immediately before the program instructions that implement
$\theta$. Let $t,S,p\in\mathbb{N}$ and let $e$ be a register record such that the
template's initial-step flag is true exactly when $t=0$, $e.q=0$ when $t=0$, and
$e.q=e.i-S$ when $t\ne 0$, with $S\le e.i$ whenever $t\ne 0$. Starting at the label for
$\theta$ with register values given by $e$, input position $p$, and any output prefix,
the counter program reaches the successor label in at most $K$ steps, leaves the input
position unchanged, and appends the word obtained by executing the template's operations
using the register values in $e$; the final register record is $e$ with $i$ replaced by
$e.i+1$ and $q$ replaced by $0$ if $t=0$ and by $e.i+1-S$ otherwise, where $K$ is the
maximum length of a template program plus three.
\end{theorem}

\begin{theorem}[Successor formula for a shifted range]
\lean{Complexity.UTab.range'_succ_left}
\label{Complexity.UTab.range'_succ_left}
\leanok
\difficulty{1}
For natural numbers $s$ and $n$, the list of the first $n+1$ consecutive natural numbers
beginning at $s$ is $[s]$ followed by the list of the first $n$ consecutive numbers
beginning at $s+1$.
\end{theorem}

\begin{theorem}[Natural-number range as a shifted list]
\lean{Complexity.UTab.range'_eq_map}
\label{Complexity.UTab.range'_eq_map}
\leanok
\difficulty{1}
For natural numbers $s,n$, the list of $s$ consecutive natural numbers starting at $n$
is the list of the first $n$ natural numbers with $s$ added to each entry.
\end{theorem}

\begin{definition}[Configuration stride]
\lean{Complexity.UTab.Str}
\label{Complexity.UTab.Str}
\leanok
\uses{Complexity.cfgStride}
For a finite binary Turing machine $M$, a bit string $z$, and natural numbers $n,T$,
define the stride to be the configuration stride for $M$ at width $|z|+n$ and time bound
$T$.
\end{definition}

\begin{theorem}[Register-copying loop from a register state]
\lean{Complexity.UTab.goes_copyR}
\label{Complexity.UTab.goes_copyR}
\leanok
\uses{Complexity.CounterProg.Goes, Complexity.UTab.CpS, Complexity.UTab.CpS.dst, Complexity.UTab.CpS.exit, Complexity.UTab.CpS.src, Complexity.UTab.ER, Complexity.UTab.ER.f, Complexity.UTab.ER.set, Complexity.UTab.ER.update_f, Complexity.UTab.goes_copy, Complexity.UTab.prog}
\difficulty{1}
Let $M$ be a finite-state multitape Turing machine with Boolean tape symbols, let $x$ be
a finite Boolean input, and let $s$ be a copy-loop mode. Given an eleven-register state
$e$ whose scratch register is zero, and an input position $p$, the program reaches the
exit mode with the destination register increased by the initial source-register value,
all other register values unchanged, and input position still $p$. For every output
prefix, this occurs within $7a+2$ steps, where $a$ is the initial value of the source
register, and the program leaves the output unchanged.
\end{theorem}

\begin{theorem}[Register clearing loop]
\lean{Complexity.UTab.goes_clearR}
\label{Complexity.UTab.goes_clearR}
\leanok
\uses{Complexity.CounterProg.Goes, Complexity.UTab.ClS, Complexity.UTab.ClS.exit, Complexity.UTab.ClS.reg, Complexity.UTab.ER, Complexity.UTab.ER.f, Complexity.UTab.ER.set, Complexity.UTab.ER.update_f, Complexity.UTab.Lb, Complexity.UTab.Lmax, Complexity.UTab.goes_clear, Complexity.UTab.prog}
\difficulty{1}
Let $M$ be a finite-state multitape Turing machine with Boolean tape symbols, let $x$ be
a finite Boolean input, and let $s$ be a clearing-loop mode, $e$ an eleven-register
state, and $p$ an input position. From the program state at the start of that loop, with
register values given by $e$, the program reaches the exit mode with the register
associated to $s$ set to zero, all other register values unchanged, and input position
still $p$. For every output prefix, this occurs within $2r+1$ steps, where $r$ is the
initial value of the register associated to $s$, and the program leaves the output
unchanged.
\end{theorem}

\begin{theorem}[Addition for consecutive natural-number ranges]
\lean{Complexity.UTab.range'_add}
\label{Complexity.UTab.range'_add}
\leanok
\difficulty{3}
For natural numbers $s,a,b$, the list of $s$ consecutive natural numbers starting at
$a+b$ is the concatenation of the list of $s$ consecutive natural numbers starting at
$a$ and the list of $b$ consecutive natural numbers starting at $s+a$.
\end{theorem}

\begin{theorem}[Simulation of a Work-Cell Template]
\lean{Complexity.UTab.goes_cell_one}
\label{Complexity.UTab.goes_cell_one}
\leanok
\uses{Complexity.CounterProg.Goes, Complexity.UTab.ER, Complexity.UTab.ER.f, Complexity.UTab.Ki, Complexity.UTab.Str, Complexity.UTab.Tpl.after, Complexity.UTab.cell_bits, Complexity.UTab.goes_instr_inv, Complexity.UTab.ibits, Complexity.UTab.prog, Complexity.UTab.tpL, Complexity.cellIdx, Complexity.cfgStride}
\difficulty{4}
Let $M$ be a finite-state Turing machine with Boolean input, and let $x,z$ be finite
Boolean words. Let $n,T,t,p,r$ be natural numbers, let $\tau$ be a work-tape index, let
$\mathrm{ph}\in\{0,1,2,3,4\}$, and let $e$ be an eleven-register assignment. Write
$S=M.k(2T+1)+(z.length+n)+3$. Assume $r<2T+1$, $t\le T$, the registers of $e$ satisfy
$e_{\mathrm{nc}}=n$, $e_z=z.length$, $e_i=tS+\tau(2T+1)+r$, $e_q=0$ when $t=0$ and
$e_q=e_i-S$ otherwise, and $e_{\mathrm{lb}}=tS$. Also assume that the truth values of
$\mathrm{ph}\ne0$, $\mathrm{ph}=2$, and $\mathrm{ph}\ne4$ agree respectively with those
of $1\le r$, $r=T$, and $r+1<2T+1$. Starting at the label for the work cell indexed by
$\tau$ and phase $\mathrm{ph}$, with time-zero flag determined by $t=0$, registers $e$,
input position $p$, and any prior output word, the counter program reaches the label
following that cell template. Its registers are those of $e$ with $i$ replaced by
$e_i+1$ and $q$ by $0$ if $t=0$, or by $e_i+1-S$ otherwise; it appends the instruction
bits for position $e_i$, in at most $L_{\max}(M)+3$ steps.
\end{theorem}

\begin{theorem}[Traversal of Work Tape Cells]
\lean{Complexity.UTab.goes_cells}
\label{Complexity.UTab.goes_cells}
\leanok
\uses{Complexity.CounterProg.Goes, Complexity.CounterProg.Goes.congr, Complexity.CounterProg.Goes.trans, Complexity.CounterProg.goes_loop, Complexity.UTab.CpS.dst, Complexity.UTab.CpS.exit, Complexity.UTab.CpS.src, Complexity.UTab.ER, Complexity.UTab.ER.f, Complexity.UTab.ER.f_ll, Complexity.UTab.ER.f_t, Complexity.UTab.ER.set_ll, Complexity.UTab.ER.update_f, Complexity.UTab.Hd, Complexity.UTab.Hd.body, Complexity.UTab.Hd.exit, Complexity.UTab.Hd.reg, Complexity.UTab.Ki, Complexity.UTab.Lb, Complexity.UTab.Lmax, Complexity.UTab.Str, Complexity.UTab.Tpl.after, Complexity.UTab.cellStart, Complexity.UTab.goes_cell_one, Complexity.UTab.goes_copyR, Complexity.UTab.goes_decR, Complexity.UTab.ibits, Complexity.UTab.prog, Complexity.UTab.rLL, Complexity.UTab.range'_add, Complexity.UTab.range'_eq_map, Complexity.UTab.tpL, Complexity.cellIdx}
\difficulty{6}
Let $M$ be a multitape Turing machine with Boolean symbols, let $x$ be a finite Boolean
input, and let $n,T,t$ be natural numbers, $z$ a finite Boolean string, $\tau$ a
work-tape index of $M$, $e$ an eleven-register record, and $p$ an input position.
Suppose $1 \le T$, $t \le T$, and the registers of $e$ satisfy $e.nc=n$, $e.z=|z|$,
$e.t=T$, $e.i=t\,\mathrm{Str}(M,n,z,T)+\tau(2T+1)$, $e.q=0$ when $t=0$ and
$e.q=e.i-\mathrm{Str}(M,n,z,T)$ otherwise, $e.lb=t\,\mathrm{Str}(M,n,z,T)$, and
$e.ll=e.tmp=0$. Starting in the program template for work-tape cell $\tau$ at phase $0$,
with register values $e.f$ and input position $p$, the program reaches the start of the
next work-tape cell; its registers are $e.f$ with $i$ replaced by $e.i+(2T+1)$ and $q$
replaced by $0$ when $t=0$ or by $e.i+(2T+1)-\mathrm{Str}(M,n,z,T)$ otherwise, while the
input position remains $p$. During this computation it appends the concatenation, over
$j=e.i,e.i+1,\ldots,e.i+2T$, of the bit strings generated by the tabulated instruction
at index $j$, and takes at most $(2T+3)(\mathrm{Lmax}(M)+20)$ steps.
\end{theorem}

\begin{theorem}[Traversal across a tape layer]
\lean{Complexity.UTab.goes_tapes}
\label{Complexity.UTab.goes_tapes}
\leanok
\uses{Complexity.CounterProg.Goes, Complexity.CounterProg.Goes.congr, Complexity.CounterProg.Goes.trans, Complexity.UTab.ER, Complexity.UTab.ER.f, Complexity.UTab.Lmax, Complexity.UTab.Str, Complexity.UTab.cellStart, Complexity.UTab.goes_cells, Complexity.UTab.goes_refl, Complexity.UTab.ibits, Complexity.UTab.prog, Complexity.UTab.range'_add, Complexity.UTab.tpL}
\difficulty{5}
Let $M$ be a finite Turing machine with $k$ work tapes, let $x$ be a finite binary
string, let $n,T,t\in\mathbb N$, and let $z$ be a finite binary string. Suppose $1\le T$
and $t\le T$. For any $d,\tau\in\mathbb N$ and register record $e$, assume $\tau+d=k$,
$e$ records $n$ as its circuit-input count, $|z|$ as its hard-wired input length, and
$T$ as its step budget, and that its current instruction index is
$t\,\mathrm{Str}(M,n,z,T)+\tau(2T+1)$, its previous-cell index is $0$ when $t=0$ and
otherwise the current index minus $\mathrm{Str}(M,n,z,T)$, its first-instruction-of-step
index is $t\,\mathrm{Str}(M,n,z,T)$, and its inner-loop and scratch registers are both
zero. Then, for every input position $p\in\mathbb N$, the counter program started at the
beginning of the cell indexed by $\tau$ reaches the input-position phase with the same
input position and register record obtained from $e$ by adding $d(2T+1)$ to the current
instruction index and updating the previous-cell index to $0$ if $t=0$, or otherwise to
the new instruction index minus $\mathrm{Str}(M,n,z,T)$. During this run it appends the
concatenation of the instruction encodings for indices from the original current index
through the next $d(2T+1)-1$ indices, and the run takes at most
$d(2T+3)(L_{\max}(M)+20)$ steps.
\end{theorem}

\begin{theorem}[Tabulated input source]
\lean{Complexity.UTab.tabLayout_getD}
\label{Complexity.UTab.tabLayout_getD}
\leanok
\uses{Complexity.tabLayout}
\difficulty{3}
For $N,k\in\mathbb{N}$ with $k<N$, the $k$-th entry of the length-$N$ list of input-bit
sources is the source for input bit $k$.
\end{theorem}

\begin{theorem}[Input Positions in the Table Layout]
\lean{Complexity.UTab.getElem?_tabLayout}
\leanok
\uses{Complexity.tabLayout}
\difficulty{1}
For natural numbers $N$ and $k$ with $k<N$, the entry at position $k$ in the list of the
first $N$ input-bit sources is the source for input bit $k$.
\end{theorem}

\begin{theorem}[Execution at an Input Position]
\lean{Complexity.UTab.goes_inp_one}
\label{Complexity.UTab.goes_inp_one}
\leanok
\uses{Complexity.CounterProg.Goes, Complexity.UTab.ER, Complexity.UTab.ER.f, Complexity.UTab.IPh, Complexity.UTab.IPh.lay, Complexity.UTab.Ki, Complexity.UTab.SSrc.Ok, Complexity.UTab.SSrc.toBitSrc, Complexity.UTab.Str, Complexity.UTab.Tpl.after, Complexity.UTab.goes_instr_inv, Complexity.UTab.ibits, Complexity.UTab.inpPhS, Complexity.UTab.inpS, Complexity.UTab.inp_bits, Complexity.UTab.layS, Complexity.UTab.prog, Complexity.UTab.tpL, Complexity.inpIdx, Complexity.snapWidth, Complexity.tabLayout}
\difficulty{4}
Let $M$ be a finite multitape Turing machine with Boolean symbols, let $x$ be a finite
Boolean input, let $n,T$ be natural numbers, let $z$ be a finite Boolean string, and let
$t
earrow T$ and $p<|z|+n+2$. Choose an input phase $\text{ph}$ and an 11-register
assignment $e$. Suppose the registers for the input count and hard-wired length equal
$n$ and $|z|$, respectively; the instruction position is the position indexed by $p$ in
layer $t$; the preceding-cell register is zero if $t=0$ and otherwise is the instruction
position minus one layer stride; and the layer-start register equals $t$ times the layer
stride. Suppose also that the input-phase source list agrees with the input-position
source list specified by whether $p$ is $0$, $1$, or $|z|+n+1$, and that the layout
source is valid and has the value of the $(p-1)$st bit of the layout of length $|z|+n$
when $1\le p\le |z|+n$, and false otherwise. Then the counter program, starting at the
label for phase $\text{ph}$ at layer $t$ with register values $e$ and input position
$ps$, reaches the successor label for that phase within $L_{\max}(M)+3$ steps, emits the
prescribed instruction bits, and ends with instruction position increased by one and
preceding-cell register equal to zero if $t=0$, or otherwise equal to the new
instruction position minus one layer stride; all other registers and the input position
are unchanged.
\end{theorem}

\begin{theorem}[Input phase of the universal simulation]
\lean{Complexity.UTab.goes_inputs}
\label{Complexity.UTab.goes_inputs}
\leanok
\uses{Complexity.CounterProg.Goes, Complexity.CounterProg.Goes.congr, Complexity.CounterProg.Goes.mono, Complexity.CounterProg.Goes.trans, Complexity.CounterProg.goes_jz_pos, Complexity.CounterProg.goes_jz_zero, Complexity.CounterProg.goes_loop, Complexity.UTab.ClS.exit, Complexity.UTab.ClS.reg, Complexity.UTab.CpS.dst, Complexity.UTab.CpS.exit, Complexity.UTab.CpS.src, Complexity.UTab.ER, Complexity.UTab.ER.f, Complexity.UTab.ER.f_ll, Complexity.UTab.ER.f_nc, Complexity.UTab.ER.f_pa, Complexity.UTab.ER.f_pb, Complexity.UTab.ER.f_z, Complexity.UTab.ER.set, Complexity.UTab.ER.set_ll, Complexity.UTab.ER.set_pa, Complexity.UTab.ER.set_pb, Complexity.UTab.ER.update_f, Complexity.UTab.Hd.body, Complexity.UTab.Hd.exit, Complexity.UTab.Hd.reg, Complexity.UTab.IPh, Complexity.UTab.IPh.lay, Complexity.UTab.Ki, Complexity.UTab.Lb, Complexity.UTab.Lmax, Complexity.UTab.SSrc.Ok, Complexity.UTab.SSrc.toBitSrc, Complexity.UTab.Str, Complexity.UTab.Tpl.after, Complexity.UTab.getElem?_tabLayout, Complexity.UTab.goes_clearR, Complexity.UTab.goes_copyR, Complexity.UTab.goes_incR, Complexity.UTab.goes_inp_one, Complexity.UTab.ibits, Complexity.UTab.inpPhS, Complexity.UTab.inpS, Complexity.UTab.layS, Complexity.UTab.prog, Complexity.UTab.rLL, Complexity.UTab.rNC, Complexity.UTab.rPA, Complexity.UTab.rPB, Complexity.UTab.rZ, Complexity.UTab.range'_add, Complexity.UTab.range'_eq_map, Complexity.UTab.tpL, Complexity.inpIdx, Complexity.snapWidth, Complexity.tabLayout}
\difficulty{7}
Let $M$ be a finite-state Turing machine with Boolean input and $k$ work tapes, let $x$
be a finite Boolean word, let $n,T,t$ be natural numbers, and let $z$ be a finite
Boolean word, with $t\le T$. Let $e$ be a register record satisfying $e.nc=n$,
$e.z=|z|$, $e.i=t\,\mathrm{Str}(M,n,z,T)+k(2T+1)$, $e.q=0$ when $t=0$ and
$e.q=e.i-\mathrm{Str}(M,n,z,T)$ otherwise, $e.lb=t\,\mathrm{Str}(M,n,z,T)$, and
$e.ll=e.pa=e.pb=e.tmp=0$. For every natural input position $ps$, the counter program
started at the input phase with label determined by $t=0$, registers $e.f$, input
position $ps$, and any output word reaches the snapshot phase with label determined by
$t=0$, registers obtained from $e.f$ by replacing $e.i$ by $e.i+|z|+n+2$ and $e.q$ by
$0$ when $t=0$ or by $e.i+|z|+n+2-\mathrm{Str}(M,n,z,T)$ otherwise, and unchanged input
position $ps$. During this computation it appends the encoded instruction bits for each
index in the interval from $e.i$ up to but excluding $e.i+|z|+n+2$, in order, and uses
at most $(|z|+n+2)(L_{\max}(M)+20)$ steps.
\end{theorem}
